% GENERATED FILE. Edit canonical lessons, then run npm run generate:books.
% canonical-source-hash: fnv1a64:c4bbabd0e67a515d
% canonical-lessons: LA-C396-promissum, LA-C396-polliceor, LA-C396-fiducia, LA-C396-sponsio, LA-C396-obstringo, LA-R396-review-389-396

\chapter{Making a Promise}
\label{ch:la-making-a-promise}
\hlchaptermodality{396}

\begin{chapteropening}
\textbf{By the end of this chapter:} \emph{I can make a promise and say that I promised something, talk about trust and a formal pledge, and say that I am bound by a favour.}

Complete the last lesson of chapter 396: I can make a promise and say that I promised something, talk about trust and a formal pledge, and say that I am bound by a favour.
\end{chapteropening}

\section[prōmissum, prōmissī]{prōmissum, prōmissī — a promise}

\label{lesson:LA-C396-promissum}



\begin{quote}
Before the new one: say the Latin for a putting off, then the Latin for expectation.
\end{quote}

\subsection*{You'll want to know: prōmissum, prōmissī}
\textbf{prōmissum, prōmissī} — \textquotedblleft{}a promise\textquotedblright{}.

\begin{quote}
Prōmissum suum servāvit.
\end{quote}

\begin{quote}
\emph{He kept his promise.}
\end{quote}

\emph{prōmissum, prōmissī} n., the participle of \emph{prōmittō}: a thing sent forward. \emph{Prōmissum servāre}, to keep a promise. English \emph{promise}.

\begin{cousinweb}
\emph{prōmissum} is the participle of \emph{prōmittō} \textquotedblleft{}to send forward, to promise\textquotedblright{} (EtymDB; Lewis \& Short), from \emph{prō} and \emph{mittō} \textquotedblleft{}to send\textquotedblright{}. English \emph{promise} comes from it through French.
\end{cousinweb}

\begin{grammarlens}[title={What you've built}]
One more word for telling, explaining and describing at B1.
\end{grammarlens}

\subsection*{Your turn}
\begin{itemize}
\raggedright
  \item \emph{Say it:} \emph{prōmissum, prōmissī}
  \item \emph{Say it:} \emph{prōmissum, prōmissī}, once more
  \item \emph{Say it:} a sentence of your own with \emph{prōmissum, prōmissī}
  \item \emph{Recall:} say the Latin for a putting off, then the Latin for expectation, then say \emph{prōmissum, prōmissī} again
\end{itemize}

\begin{tcolorbox}[breakable,colback=teal!4,colframe=teal!35!black,title={Before you move on}]
What does prōmissum, prōmissī mean? (\textquotedblleft{}A promise\textquotedblright{}.) Say it once more.
\end{tcolorbox}
\section[polliceor, pollicērī]{polliceor, pollicērī — to promise, to offer}

\label{lesson:LA-C396-polliceor}



\begin{quote}
Before the new one: say the Latin for expectation, then the Latin for a promise.
\end{quote}

\subsection*{You'll want to know: polliceor, pollicērī}
\textbf{polliceor, pollicērī} — \textquotedblleft{}to promise, to offer\textquotedblright{}.

\begin{quote}
Pater mihi equum pollicitus est.
\end{quote}

\begin{quote}
\emph{Father promised me a horse.}
\end{quote}

Deponent: \emph{polliceor, pollicērī, pollicitus sum}. Often with the accusative and future infinitive: \emph{pollicētur sē ventūrum esse}, he promises that he will come.

\begin{cousinweb}
\emph{polliceor} is built from an old prefix \emph{por-} \textquotedblleft{}forth\textquotedblright{} and \emph{liceor} \textquotedblleft{}to bid, to offer a price\textquotedblright{} (EtymDB; Lewis \& Short): to hold out an offer.
\end{cousinweb}

\begin{grammarlens}[title={What you've built}]
One more word for telling, explaining and describing at B1.
\end{grammarlens}

\subsection*{Your turn}
\begin{itemize}
\raggedright
  \item \emph{Say it:} \emph{polliceor, pollicērī}
  \item \emph{Say it:} \emph{polliceor, pollicērī}, once more
  \item \emph{Say it:} a sentence of your own with \emph{polliceor, pollicērī}
  \item \emph{Recall:} say the Latin for expectation, then the Latin for a promise, then say \emph{polliceor, pollicērī} again
\end{itemize}

\begin{tcolorbox}[breakable,colback=teal!4,colframe=teal!35!black,title={Before you move on}]
What does polliceor, pollicērī mean? (\textquotedblleft{}To promise, to offer\textquotedblright{}.) Say it once more.
\end{tcolorbox}
\section[fīdūcia, fīdūciae]{fīdūcia, fīdūciae — trust, confidence; a thing held in trust}

\label{lesson:LA-C396-fiducia}



\begin{quote}
Before the new one: say the Latin for a promise, then the Latin for to promise.
\end{quote}

\subsection*{You'll want to know: fīdūcia, fīdūciae}
\textbf{fīdūcia, fīdūciae} — \textquotedblleft{}trust, confidence; a thing held in trust\textquotedblright{}.

\begin{quote}
Magnam fīdūciam in amīcīs habēbat.
\end{quote}

\begin{quote}
\emph{He had great trust in his friends.}
\end{quote}

\emph{fīdūcia, fīdūciae} f., built on \emph{fīdō}, to trust. In law it is a thing handed over in trust. English \emph{fiduciary}.

\begin{cousinweb}
\emph{fīdūcia} is built on \emph{fīdō} \textquotedblleft{}to trust\textquotedblright{} (EtymDB; Lewis \& Short), the verb behind \emph{fīdus} \textquotedblleft{}faithful\textquotedblright{}. English \emph{fiduciary} comes from it.
\end{cousinweb}

\begin{grammarlens}[title={What you've built}]
One more word for telling, explaining and describing at B1.
\end{grammarlens}

\subsection*{Your turn}
\begin{itemize}
\raggedright
  \item \emph{Say it:} \emph{fīdūcia, fīdūciae}
  \item \emph{Say it:} \emph{fīdūcia, fīdūciae}, once more
  \item \emph{Say it:} a sentence of your own with \emph{fīdūcia, fīdūciae}
  \item \emph{Recall:} say the Latin for a promise, then the Latin for to promise, then say \emph{fīdūcia, fīdūciae} again
\end{itemize}

\begin{tcolorbox}[breakable,colback=teal!4,colframe=teal!35!black,title={Before you move on}]
What does fīdūcia, fīdūciae mean? (\textquotedblleft{}Trust, confidence; a thing held in trust\textquotedblright{}.) Say it once more.
\end{tcolorbox}
\section[spōnsiō, spōnsiōnis]{spōnsiō, spōnsiōnis — a solemn promise, a pledge; a wager}

\label{lesson:LA-C396-sponsio}



\begin{quote}
Before the new one: say the Latin for to promise, then the Latin for trust.
\end{quote}

\subsection*{You'll want to know: spōnsiō, spōnsiōnis}
\textbf{spōnsiō, spōnsiōnis} — \textquotedblleft{}a solemn promise, a pledge; a wager\textquotedblright{}.

\begin{quote}
Spōnsiōnem inter sē fēcērunt.
\end{quote}

\begin{quote}
\emph{They made a formal wager with each other.}
\end{quote}

\emph{spōnsiō, spōnsiōnis} f., built on \emph{spondeō}, to pledge: the old spoken contract, \emph{spondēsne?} — \emph{spondeō}. Also a bet. English \emph{sponsor} is from the same verb.

\begin{cousinweb}
\emph{spōnsiō} is the noun built on \emph{spondeō} \textquotedblleft{}to pledge solemnly\textquotedblright{} (EtymDB; Lewis \& Short). English \emph{sponsor} comes from the same verb.
\end{cousinweb}

\begin{grammarlens}[title={What you've built}]
One more word for telling, explaining and describing at B1.
\end{grammarlens}

\subsection*{Your turn}
\begin{itemize}
\raggedright
  \item \emph{Say it:} \emph{spōnsiō, spōnsiōnis}
  \item \emph{Say it:} \emph{spōnsiō, spōnsiōnis}, once more
  \item \emph{Say it:} a sentence of your own with \emph{spōnsiō, spōnsiōnis}
  \item \emph{Recall:} say the Latin for to promise, then the Latin for trust, then say \emph{spōnsiō, spōnsiōnis} again
\end{itemize}

\begin{tcolorbox}[breakable,colback=teal!4,colframe=teal!35!black,title={Before you move on}]
What does spōnsiō, spōnsiōnis mean? (\textquotedblleft{}A solemn promise, a pledge; a wager\textquotedblright{}.) Say it once more.
\end{tcolorbox}
\section[obstringō, obstringere]{obstringō, obstringere — to bind, to put under an obligation}

\label{lesson:LA-C396-obstringo}



\begin{quote}
Before the new one: say the Latin for trust, then the Latin for a solemn promise.
\end{quote}

\subsection*{You'll want to know: obstringō, obstringere}
\textbf{obstringō, obstringere} — \textquotedblleft{}to bind, to put under an obligation\textquotedblright{}.

\begin{quote}
Beneficiō mē obstrīnxistī.
\end{quote}

\begin{quote}
\emph{You have bound me by your kindness.}
\end{quote}

\emph{obstringō, obstringere, obstrīnxī, obstrictum}: to tie fast, then to bind by a promise or a favour.

\begin{cousinweb}
\emph{obstringō} is built from \emph{ob} and \emph{stringō} \textquotedblleft{}to draw tight, to bind\textquotedblright{} (Lewis \& Short): to tie someone fast to a duty.
\end{cousinweb}

\begin{grammarlens}[title={What you've built}]
One more word for telling, explaining and describing at B1.
\end{grammarlens}

\subsection*{Your turn}
\begin{itemize}
\raggedright
  \item \emph{Say it:} \emph{obstringō, obstringere}
  \item \emph{Say it:} \emph{obstringō, obstringere}, once more
  \item \emph{Say it:} a sentence of your own with \emph{obstringō, obstringere}
  \item \emph{Recall:} say the Latin for trust, then the Latin for a solemn promise, then say \emph{obstringō, obstringere} again
\end{itemize}

\begin{tcolorbox}[breakable,colback=teal!4,colframe=teal!35!black,title={Before you move on}]
What does obstringō, obstringere mean? (\textquotedblleft{}To bind, to put under an obligation\textquotedblright{}.) Say it once more.
\end{tcolorbox}
\section[(dialogue)]{Words from chapters 389–396 — a review}

\label{lesson:LA-R396-review-389-396}



\begin{quote}
Say the Latin for slavery and to set free.
\end{quote}

\subsection*{Your turn}
Say these aloud:

\begin{itemize}
\raggedright
  \item slavery, to set free, dutifulness, a load, compliance
  \item a condition, but if, provided that, only, however much
  \item even if, although, otherwise, differently, or if
  \item to beware, caution, necessity, convenient, to take precautions
  \item to set up, to fix, a plan, an undertaking, to try
  \item to prepare, to set out, to set in order, to set about, to foresee
  \item to hasten, to put off, to extend, a putting off, expectation
  \item a promise, to promise, trust, a solemn promise, to bind
\end{itemize}

\begin{tcolorbox}[breakable,colback=teal!4,colframe=teal!35!black,title={Before you move on}]
Slavery? (\textbf{servitūs, servitūtis}.) To set free? (\textbf{manūmittō, manūmittere}.) Dutifulness? (\textbf{pietās, pietātis}.) A load? (\textbf{onus, oneris}.) Compliance? (\textbf{obsequium, obsequiī}.) A condition? (\textbf{condiciō, condiciōnis}.) But if? (\textbf{sīn}.) Provided that? (\textbf{dummodo}.) Only? (\textbf{modo}.) However much? (\textbf{quamvīs}.) Even if? (\textbf{etiamsī}.) Although? (\textbf{tametsī}.) Otherwise? (\textbf{aliōquī}.) Differently? (\textbf{secus}.) Or if? (\textbf{sīve}.) To beware? (\textbf{caveō, cavēre}.) Caution? (\textbf{cautiō, cautiōnis}.) Necessity? (\textbf{necessitās, necessitātis}.) Convenient? (\textbf{opportūnus, opportūna, opportūnum}.) To take precautions? (\textbf{praecaveō, praecavēre}.) To set up? (\textbf{statuō, statuere}.) To fix? (\textbf{dēstinō, dēstināre}.) A plan? (\textbf{prōpositum, prōpositī}.) An undertaking? (\textbf{inceptum, inceptī}.) To try? (\textbf{cōnor, cōnārī}.) To prepare? (\textbf{apparō, apparāre}.) To set out? (\textbf{dispōnō, dispōnere}.) To set in order? (\textbf{ōrdinō, ōrdināre}.) To set about? (\textbf{mōlior, mōlīrī}.) To foresee? (\textbf{praevideō, praevidēre}.) To hasten? (\textbf{mātūrō, mātūrāre}.) To put off? (\textbf{differō, differre}.) To extend? (\textbf{prōrogō, prōrogāre}.) A putting off? (\textbf{dīlātiō, dīlātiōnis}.) Expectation? (\textbf{exspectātiō, exspectātiōnis}.) A promise? (\textbf{prōmissum, prōmissī}.) To promise? (\textbf{polliceor, pollicērī}.) Trust? (\textbf{fīdūcia, fīdūciae}.) A solemn promise? (\textbf{spōnsiō, spōnsiōnis}.) To bind? (\textbf{obstringō, obstringere}.)
\end{tcolorbox}
